\subsection{From Thermoelastic Formulation to Prediction Task}

The mathematical formulation in Chapter 2 describes thermoelastic wave propagation through a set of coupled physical fields and material parameters. In its complete form, the problem involves the temperature field \(T(x,t)\), the temperature perturbation \(\theta(x,t)\), the displacement field \(u(x,t)\), the strain tensor \(\varepsilon_{ij}\), the stress tensor \(\sigma_{ij}\), and the material parameters of the geological medium. These quantities are connected through the equation of motion, the strain--displacement relation, the thermoelastic constitutive relation, and the coupled heat-conduction equation.

For the purposes of this thesis, this formulation is interpreted as the physical basis of a prediction task. Instead of claiming to solve the complete coupled partial differential equation system for every point in a two-dimensional field-scale geological domain, the methodology defines a structured mapping from a set of physically meaningful inputs to a set of predicted thermoelastic response quantities. In general form, this can be written as
\begin{equation}
\mathcal{F}:
X
\rightarrow
\widehat{Y},
\end{equation}
where \(X\) denotes the input description of the material, thermal state, spatial configuration, and propagation direction, while \(\widehat{Y}\) denotes the predicted response.

The input description \(X\) is based on the variables introduced in Chapter 2. It may include material type, density, elastic parameters, thermal parameters, temperature or temperature perturbation, source location, probe location, propagation direction, and time or observation settings. These quantities provide a compact representation of the thermoelastic problem without requiring the full spatial field to be explicitly computed everywhere in the domain.

The predicted response \(\widehat{Y}\) represents the estimated behavior of the thermoelastic disturbance along a selected direction or at a selected observation point. Depending on the model component and available output structure, this response may include predicted temperature response, displacement components, directional wave behavior, or comparative model-level indicators. The purpose is not to treat these outputs as direct field-validated measurements, but to use them as structured estimates that can be compared across materials, propagation directions, and model families.

This prediction-oriented formulation is useful because directional behavior is central to geological wave propagation. In an ideal homogeneous and isotropic medium, the material response would not depend on the direction of propagation. In real geological media, bedding, fractures, aligned pores, lithological contrasts, and anisotropic thermal or elastic properties may cause the response to vary with direction. The source--probe geometry introduced in Chapter 2 therefore becomes a key part of the methodology.

The conversion from a physical formulation to a prediction task can be summarized as follows. Chapter 2 defines the physical variables and governing relations. The present chapter selects the variables that can be used as inputs and outputs in a predictive framework. The resulting methodology allows different predictive components to be compared under the same physical assumptions and with the same interpretation of material parameters and propagation direction.
